The compilation proceeds by analyzing the structure of the instance type.
There are three relevant cases. Rules \eI{~\PYG{c}{r1}~} and
\eI{~\PYG{c}{r2}~} handle quantifications, while rule
\eI{~\PYG{c}{r3}~} is the base case that constructs the final \elpi rule
from the information accumulated during the recursive traversal.

When the current type has the form \eI{(prod _ Ty Bo)}, the abstraction binds
a variable of type \eI{Ty} in the body \eI{Bo}. The test \eI{is-class? Ty}
checks whether \eI{Ty} has the shape
$\forall x_1 \dots x_n, T$ (possibly with zero quantifiers), where \eI{T} is a
term whose applicative head is a type class.

If this condition holds, we are compiling a conditional instance and rule
\eI{~\PYG{c}{r1}~} is applied. The rule introduces a fresh nominal variable
\eI{x} representing both the abstracted variable in \eI{Bo} and the local
instance with type \eI{Ty}. To construct the premise
corresponding to \eI{x}, the compiler recursively invokes \eI{comp} on \eI{x}
while switching the polarity of the rule. This recursive call produces the
premise \eI{M}, which abstracts \eI{x}. Compilation then continues on
\eI{Bo x}, while extending the accumulators with the new instance argument
\eI{x} and the new premise \eI{M x}.

Rule \eI{~\PYG{c}{r2}~} handles the case where the quantified type is not a
type class. In this situation, the compiler simply consumes under the
abstraction, leaving the list of premises \eI{P} unchanged and adding the
variable \eI{x} to the list of instance arguments.

Both rules \eI{~\PYG{c}{r1}~} and \eI{~\PYG{c}{r2}~} return a term of the form
\eI{(pi y\ R y)}, where \eI{R} is the result of the recursive call to
\eI{comp}. Since \eI{R} has type \eI{term -> term} and abstracts the local
variable \eI{x}, the compiler introduces a fresh variable \eI{y} to quantify
over the corresponding argument. This step ensures that the resulting rule is
closed (ground), which is required by the \elpi language.

The base case of \eI{comp} is reached when the current type \eI{Ty} is not a
quantification. In this case, \eI{Ty} is a type class \eI{C} applied to a list
of arguments \eI{CAg}. The compiler then constructs the final rule by building
its head and premises.

To build the head \eI{Head}, the compiler first constructs the proof term
\eI{Proof}, obtained by applying the instance \eI{I} to its arguments
\eI{rev Ag}. The list \eI{Ag} is stored in reverse order during traversal, so
it is reversed before constructing the application. The predicate
\eI{make-rule-head} takes the class name as its first argument in order to
generate the predicate corresponding to the class. The arguments of this
predicate are obtained by appending the proof term to the list of class
arguments.

The premises of the clause are stored in the accumulator \eI{P}. As with the
instance arguments, they are reversed before constructing the final rule.

Finally, the predicate \eI{build-rule} assembles the rule according to the
current polarity. With positive polarity, it produces a standard clause
\eI{Head :- Prems}. With negative polarity, it produces the implication
\eI{Prems => Head}.
